\chapter{Design som verdsskaping}

\begin{motto}
Vi formar verda, og verda formar oss attende.\\ Det er ikkje ein metafor. Det er fundamentet.
\end{motto}

Heile saka mot designfaget har ein botn, og botnen er økonomisk. Design, heiter det til sist, er eit underbruk av marknaden — ein måte å gjere ting sal-bare på, å pakke vara, å pynte transaksjonen. Faget kan ikkje seie nei til oppdragsgjevaren, og difor har det inga eiga grunn å stå på; det er ein tenar utan eige hus. Dette kapitlet hevdar det stikk motsette, og det er det djupaste forsvaret boka har å føre. Design er ikkje eit underbruk av økonomien. Design er \emph{ontologisk} — det er verdsskapande — og denne innsikta gjev faget eit fundament djupare enn noko marknaden kan gje eller ta. Når vi formar tinga kring oss, formar vi samstundes måtane å leve på som tinga gjer moglege, og dermed formar vi oss sjølve. Eit fag som gjer dette, svarar ikkje fyrst til marknaden. Det svarar til verda det er med på å lage, og til dei som skal bu i henne.

\section{Det ontologiske grepet}

Den klaraste forma for innsikta kom frå Anne-Marie Willis, som gav henne namnet \emph{ontological designing} \autocite{willis2006ontological}. Tesen hennar lèt seg seie i éi setning, og ho har vorte ståande som ei læresetning for heile retninga: vi designar verda vår, medan verda vår verkar attende og designar oss. Det er ein sirkel, og sirkelen er poenget. Vi lagar ting — ein by, ein skrift, eit grensesnitt — og desse tinga lagar i sin tur vanane, rørslene, samværsformene og sjølvforståinga til dei som lever mellom dei. Motorvegen designar ikkje berre reisa; han designar byen, nærleiken, kva eit nabolag er. Skjermen designar ikkje berre lesinga; han designar merksemda, samtalen, kva det vil seie å vere åleine eller saman. Design er såleis aldri berre å lage gjenstandar. Det er å gripe inn i sjølve det å vere menneske i verda.

Dette er ikkje ein høgtfløygen påstand utan rot; det er ei filosofisk innsikt med tyngd, henta frå fenomenologien sitt syn på at mennesket og verda blir til i lag, og ført over på praksisen som faktisk lagar den verda. Og det er nett her fundamentet ligg. For om design er verdsskapande, då er ikkje spørsmålet om faget «sel» eller «pyntar». Spørsmålet er kva slags verd det er med på å lage, og kven den verda tener. Den som reduserer design til marknadstenest, har ikkje sett gjennom faget. Han har sett bort frå det djupaste det gjer. Designaren formar ikkje berre eit produkt; han formar, i det små og i det store, korleis det skal vere å vere til. Eit ansvar av det slaget er ikkje fråvêret av eit fundament. Det er eit fundament så tungt at faget framleis strevar med å bere det.

\section{Å lage framtid, å øydeleggje henne}

Tony Fry gjorde det ontologiske grepet til ein etikk, og det er hos Fry at saka mot den reine marknadslogikken blir skarpast. I \emph{Design Futuring} \autocite{fry2009} la han fram eit par omgrep som snur synet på kva design eigentleg gjer: \emph{defuturing} og \emph{futuring}. Defuturing er det design gjer når det tek framtid frå oss — når tinga vi lagar, og levesettet dei ber, et opp dei vilkåra som gjer framtidig liv mogleg. Mykje av det moderne designet, hevdar Fry, har vore defuturande utan å vite det: kvar veldesigna eingongsting, kvar smart løysing som bind oss djupare til eit ikkje-haldbart system, tek litt framtid frå dei som kjem etter. Mot dette set han futuring: design som medvite lagar vilkår for at det skal finnast ei framtid å leve i. Skiljet er ikkje smak og ikkje sal. Det er etisk, og det er ontologisk, fordi det handlar om kva slags verd og kva slags tid tinga våre opnar eller stengjer.

Det avgjerande hos Fry er at design her får ein \emph{eigen} målestokk, uavhengig av oppdragsgjevaren. Eit produkt kan vere lønsamt og likevel defuturande; det kan glede kunden og likevel stele framtid. Då er marknaden sin dom og faget sin dom ikkje den same dommen, og faget har sin eigen å felle. Dette er det presise motsvaret til skuldinga om at design ikkje kan seie nei. Fry viser kva designet skulle seie nei \emph{til}, og kvifor — ikkje av smak, men fordi den ontologiske krafta i faget gjer det medansvarleg for verda det lagar. Eit fag som kan grunngje ein nei-dom utan å låne grunngjevinga frå marknaden, har funne noko å stå på.

\section{Mot den økonomistiske moderniteten}

Arturo Escobar tok innsikta heilt ut og snudde ho mot moderniteten sjølv. I \emph{Designs for the Pluriverse} \autocite{escobar2018pluriverse} argumenterer han for at den dominerande forma for design — den som tener marknad, vekst og éin universell måte å vere moderne på — er sjølve problemet, ikkje løysinga. Han kallar denne forma for tener av \emph{the one-world world}: tanken om at det finst éi rett verd, den vestlege og økonomiske, som alle andre verder skal smelte inn i. Mot dette set han \emph{pluriversen} — at mange verder kan finnast saman — og ein \emph{autonom design} som spring ut av lokalsamfunn sjølve, retta mot rettferd, mot sjølvråderett, og mot jorda snarare enn mot vekst. Her er design ikkje lenger ein tenar for økonomien. Det er ei kraft som \emph{utfordrar} den økonomistiske moderniteten frå innsida, ved å lage andre verder enn den marknaden vil ha.

Poenget for denne boka er ikkje å avgjere om Escobar har rett i alt. Poenget er at sjølve eksistensen av eit slikt prosjekt fellar skuldinga om at design er underordna marknaden av naudsyn. Tvert om: her er ein heil designtradisjon som tek den ontologiske innsikta — at vi lagar verder — og brukar henne til å gjere motstand mot den eine verda kapitalen vil ha. Victor Margolin peika i same lei med samlinga \emph{The Politics of the Artificial} \autocite{margolin2002politics}, der han argumenterte for at den menneskeskapte verda er eit politisk felt og at design difor uunngåeleg er politisk handling, ikkje nøytral teneste. Margolin ville at faget skulle ta ansvar for «det kunstige» som heilskap — for kva slags materiell verd vi byggjer — og det er eit ansvar som ligg langt over og forut for kvart einskild oppdrag. Når faget kan reise eit slikt krav til seg sjølv, er det ikkje grunnlaust. Det har ein grunn, og grunnen er at det formar verda.

\section{Frå kritikk til oppbygging}

Ein kunne innvende at alt dette er kritikk, og at kritikk er lett: det er enkelt å seie kva design ikkje skal vere. Men den ontologiske retninga er ikkje berre kritisk; ho er òg oppbyggjande, og det er denne dobbeltheita som gjer henne til eit fundament snarare enn ein protest. Ezio Manzini har vist den oppbyggjande sida i \emph{Design, When Everybody Designs} \autocite{manzini2015}, ein heil ramme for design som driv fram sosial innovasjon — der designaren ikkje lagar ferdige løysingar for passive mottakarar, men legg til rette for at folk sjølve kan finne nye måtar å bu, ete, dele og leve saman på. Her blir den ontologiske innsikta praktisk: sidan design lagar levesett, kan det medvite lage \emph{betre} levesett, meir samhaldande og meir berekraftige. Det er verdsskaping vendt mot det gode, og det skjer i konkrete prosjekt, ikkje berre i teori.

På den andre kanten står det kritiske og spekulative designet, som brukar faget til å tenkje snarare enn å selje. Anthony Dunne og Fiona Raby gav retninga form i \emph{Speculative Everything} \autocite{dunne2013}: design som lagar ting og scenario ikkje for marknaden, men for å stille spørsmål — om kva slags framtider vi er på veg mot, om kva vi tek for gjeve, om kva som kunne vore annleis. Eit spekulativt objekt sel ingenting; det opnar ein samtale om verda. Og Carl DiSalvo viste i \emph{Adversarial Design} \autocite{disalvo2012} at design kan vere ein form for politisk usemje i seg sjølv — at det kan lage ting som gjer konflikt og maktstrid synleg og diskuterbar, i staden for å glatte over han. Begge desse er bevis, i praksis, på at design kan stå \emph{utanfor} marknadslogikken og likevel vere fullverdig design. Eit fag som kan kritisere, spekulere og gjere motstand, har openbert ikkje marknaden som einaste grunn.

Sett saman teiknar desse røystene eit fundament den økonomiske skuldinga ikkje når. Willis gjev innsikta: vi lagar verder, og verdene lagar oss. Fry gjer henne til etikk; Escobar og Margolin gjer henne til politikk; Manzini gjer henne til sosialt oppbyggjande arbeid; Dunne, Raby og DiSalvo gjer henne til kritisk tanke. Ingen av desse er underbruk av økonomien. Alle tek utgangspunkt i at design har ei eiga, verdsskapande kraft, og at krafta ber eit ansvar marknaden korkje gav eller kan oppheve. Dette er ikkje eit fag på leit etter ein grunn. Det er eit fag som har funne den djupaste grunnen av alle — at det er med på å lage verda — og som strevar, ærleg og ufullenda, med å leve opp til det.

\vignett

Det står att eitt spørsmål: om faget verkeleg har alt dette — kunnskap, metode, meining, ontologisk tyngd — kvifor er det då framleis i strid med seg sjølv om sitt eige grunnlag? Det er der dei to bøkene møtest.

\vidarelesing{\fullcite{willis2006ontological}; \fullcite{fry2009}; \fullcite{escobar2018pluriverse}; \fullcite{margolin2002politics}; \fullcite{manzini2015}; \fullcite{dunne2013}; \fullcite{disalvo2012}.}
